% -*- coding: utf-8; mode: latex; -*-

% +++
% latex = "lualatex"
% +++

%
% FIT2024 向け LaTeX クラスファイル
% https://github.com/trueroad/FITpaper-class
%
% サンプルファイル
%
% Copyright (C) 2018, 2019, 2022-2024 Masamichi Hosoda.
% All rights reserved.
%

\documentclass[dvipdfmx]{FITpaper}

% 図の貼り込み用
\usepackage{graphicx}

\usepackage{dblfloatfix}

\usepackage{multirow}

\usepackage{booktabs}

% 最終ページで両カラムの下端を揃える
% \usepackage[balance]{nidanfloat}
\usepackage{flushend}

\makeatletter
\renewcommand{\ackbibsection}[1]{%
  \par
  \vspace{6bp}%
  {\centering\normalsize\sffamily #1\par}%
  \vspace{2pt}%
}
\makeatother
% 和文タイトル
\jtitle{空き家マッチングにおける借り手の熱意を代弁する\\交渉対話エージェントの試作}

% 欧文タイトル
\etitle{Prototyping a Negotiation Dialogue Agent as a Proxy for Tenant Enthusiasm in Vacant House Matching}

% 著者数：著者の数だけ c を書く
\authors{cccc}

% 和文著者名：著者名間に & を書く
% 所属番号を \affmark でつける
\jauthors{%
  福富 隆大\affmark{1} &
  長澤 史記\affmark{1} &
  白松 俊  \affmark{1} &
  XX XX   \affmark{1}
}
% 欧文著者名：著者名間に & を書く
\eauthors{%
  XXXX XXXX &
  XXXX XXXX &
  XXXX XXXX &
  XXXX XXXX
}

% 所属
% \affmark でつけた番号毎に指定
\afftext{1}{名古屋工業大学  Nagoya Institute of Technology.}


\begin{document}

\maketitle

% ======================================
\section{はじめに}
% ======================================

日本では総務省調査（2023年）によると空き家数が約900万戸（過去最多）に達しており，
活用が求められているが，多くのオーナーは使い方や借り手の人柄次第で貸してもよいと考える消極的な層である．

このような消極的なオーナー層にはたらきかけるマッチングサービスに「さかさま不動産」がある．
従来の不動産取引では大家が借り手を選ぶのに対し，さかさま不動産では借り手が自身の活動目的や
想いを記事として公開し，大家が「この人に貸したい」と思う相手を選ぶ．
借り手は自身の夢・目的・背景を記事として公開するが，大家は対話時点で記事を詳しく読んでいない場合も多く，借り手が自ら熱意をアピールする必要がある．

こうした「情報の非対称性のもとで申請者が熱意と適合性の両方を伝える対話」は，
就職活動や研究資金申請など様々な場面に共通する課題構造であり，
多数の相手に対して自ら行うことは負担が大きく、不慣れな利用者には難しい．

本研究は，こうした状況において申請者（借り手）の記事をもとに交渉対話を代行する
AI エージェントを試作することを最終的な目標とする．
これにより，不慣れな利用者でも大家に熱意を届けられる機会を提供するとともに，
感情表現が対話品質に与える影響を実証的に明らかにする価値がある．
そのアプローチとして，感情表現の指示をプロンプトに追加する条件と追加しない条件を比較する
予備実験を行い，システムの実現可能性と評価手法を検討する．
\\

% ======================================
\section{関連研究}
% ======================================

\subsection{ロールプレイング言語エージェント（RPLA）}

LLM を用いてペルソナを与え，特定の人物として対話させるロールプレイング言語エージェント
（RPLA）の研究が進んでいる~\cite{chen2024rpla,tseng2024twotales}．
RPLA における課題として，LLM がキャラクターとして与えられた情報の範囲を超えた
事実を発話する hallucination が挙げられている~\cite{yu2024timechara}．
本研究が対象とする借り手記事は，特定個人の目的・価値観を記述した
Individualized Persona~\cite{chen2024rpla}に近く，
記事にない事実を発話するリスクへの対処として，捏造防止ルールをプロンプトに組み込んだ．

\subsection{LLMエージェントにおける感情表現}

LLM エージェントの感情表現に関して，
進化的強化学習により感情ポリシーを動的に最適化する EvoEmo~\cite{evoemo2025}，
感情状態を記憶検索に組み込む Emotional RAG~\cite{emotionalrag2024}，
戦略的計画と表現最適化を組み合わせた DMNA~\cite{dmna2025} などが提案されている．
これらの研究は，感情表現が LLM エージェントの説得力やペルソナ整合性に寄与することを示している．


% ======================================
\section{提案システム}
% ======================================

\subsection{システム概要}

提案システムは，借り手の記事を入力として大家との多ターン対話を行う
借り手エージェントである．
借り手エージェントは主に2つの機能を担う．

\begin{enumerate}
  \item \textbf{感情表現（主）}：借り手の動機を率直に伝え「この人に貸したい」と感じさせる発話を生成する．
  \item \textbf{適合確認の質問（従）}：物件が借り手の目的に合うか確認する質問を行う（2〜3回程度）．
\end{enumerate}

感情表現の効果を定量的に評価することが本研究の主目的である．

\subsection{借り手エージェントの設計}

エージェントには借り手記事の全文と発話指示をシステムプロンプトとして与える．
2条件の違いは発話指示の内容のみであり（表\ref{tab:condition}），
それ以外の借り手情報・捏造防止ルールは両条件で同一とした．

\begin{table}[t]
  \centering
  \caption{条件の比較}
  \label{tab:condition}

  \begin{tabular}{lp{4.5cm}}
    \toprule
    条件 & 発話指示の内容 \\
    \midrule
    Baseline & 借り手の目的・希望を事実に基づき伝える．感情表現指示なし． \\
    Proposed & Baseline に加え，フェーズ別の感情表現指示を付与．opening では動機を率直に，middle では懸念を受け止めてから質問，closing では前向きな姿勢を示す． \\
    \bottomrule
  \end{tabular}
\end{table}

捏造防止のため，記事に根拠のない断定的表現を避けるルールを
両条件共通でプロンプトに組み込んだ．
\\

% ======================================
\section{実験}
% ======================================

\subsection{実験設定}

提案システムの2条件（Baseline・Proposed）を ChatGPT の Custom GPTs として実装し，
人間の参加者が大家役を担う対話実験を行った．

互いに異なる借り手ケースを割り当てた6種類のGPTsを用意した（表\ref{tab:gpts}）．
参加者（研究室メンバー6名）は物件情報シートを参照しながら大家役として各GPTと4ターン程度対話し，
全6 GPTsを評価した．

\begin{table}[t]
  \centering
  \caption{実験に使用した借り手エージェント GPTs}
  \label{tab:gpts}
  \small
  \begin{tabular}{llp{1.55cm}}
    \toprule
    GPT名 & 借り手 & 条件 \\
    \midrule
    整体院 & パーソナルトレーナー・整体師 & Baseline \\
    託児カフェ & デザイナー・ベーグル屋 & Baseline \\
    陶作家 & 陶作家 & Baseline \\
    アウトドア & キャンプインストラクター & Proposed \\
    ピラティス & ピラティス講師・管理栄養士 & Proposed \\
    学童クラブ & 学童クラブ運営者 & Proposed \\
    \bottomrule
  \end{tabular}
\end{table}

\subsection{評価軸と評価方法}

表\ref{tab:survey}にアンケート項目を示す．リカートスケールは7点を最良・1点を最低とする．\\

\begin{table}[h]
  \centering
  \caption{アンケート項目}
  \label{tab:survey}
  \begin{tabular}{lp{3.2cm}l}
    \toprule
    項目 & 質問文（要約） & 尺度 \\
    \midrule
    根拠忠実性 & 記事にない事実を言っていた件数 & 違反件数 \\
    熱意の伝達 & 借り手の想いが伝わった & 7件法 \\
    人となり & どんな人かわかった & 7件法 \\
    使い方 & この物件をどう使うか想像できた & 7件法 \\
    貸したいか & 次の話に進みたいと思う & 7件法 \\
    \bottomrule
  \end{tabular}
\end{table}

% ======================================
\section{結果と考察}
% ======================================

\subsection{条件間の比較}

根拠忠実性の違反は全36件で0件であった．
条件別評価平均を表\ref{tab:results_cond}に示す．

\begin{table}[h]
  \centering
  \caption{条件別評価平均（Mann-Whitney U, $n=18$/条件）}
  \label{tab:results_cond}
  \begin{tabular}{p{2.0cm}cccc}
    \toprule
    条件 & 熱意 & 人となり & 使い方 & 貸したい \\
    \midrule
    Baseline (B) & 5.00 & 4.56 & \textbf{6.11} & 5.00 \\
    Proposed (P) & \textbf{5.50} & 4.28 & 5.17 & 5.00 \\
    \midrule
    $p$値 & .143 & .807 & \textbf{.002} & .629 \\
    \bottomrule
  \end{tabular}
\end{table}

\subsection{考察}

Mann-Whitney U 検定（両側，$n=18$/条件）の結果，
使い方のイメージスコアで Baseline が Proposed を有意に上回った
（$U=253.5, p=.002$）．
感情表現指示の付加により感情的訴求に重心が移り，
物件利用イメージの提示が相対的に少なくなった結果と考えられる．

熱意の伝達は Proposed がやや上回ったが（5.50 vs 5.00，$U=118.5, p=.143$），
人となり・貸したいを含む他の3項目でも有意差はなく，評価者数・GPTs数の拡充が必要である．

各評価者が複数 GPT を評価しているため完全独立ではなく，また各条件に異なる借り手ケースを割り当てたため（需要特性バイアス回避），条件差はケース特性との交絡を含む参考値である．

定性的コメントとして「AIらしい文章」「話が長い」「話題がずれる」などが得られ，
GPTs の単一プロンプト制約やターン数制限に起因すると考えられる．\\


% ======================================
\section{おわりに}
% ======================================

本研究では，借り手の熱意を代弁して大家と交渉対話を行う AI エージェントを試作し，
感情表現指示の有無を独立変数とした比較実験を行った．
結果，根拠忠実性の違反は0件であり，
使い方のイメージスコアでは Baseline が Proposed を有意に上回った（$p=.002$）．
熱意の伝達では Proposed がやや上回ったが有意差には至らなかった（$p=.143$）．

今後の課題として，交絡解消（被験者間カウンターバランス）と評価者数の拡大が挙げられる．

% ======================================
\acknowledgment{%
  本研究の一部はJSPS科研費（24K03052, 25K00866）の支援を受けた．
}
% ======================================

% ======================================

\footnotesize
\bibliographystyle{junsrt}
\bibliography{references}

\end{document}
